\hmsubsection{Inverse Kinematics}{UMA}{}

This section describes the inverse kinematics solver responsible for 
converting Cartesian target positions in the workspace into the 
corresponding motor positions for the arm's five Dynamixel servos. 
The solver is the geometric core of the system: when the vision 
pipeline produces the position of a detected ball in centimetres, 
the inverse kinematics solver computes the joint angles required to 
place the gripper at that position with the correct orientation, 
and converts those angles into Dynamixel step values that the 
firmware can transmit to the motors.

The solver is implemented in the \texttt{ArmIK} class in 
\texttt{src/ik/solver.py} and is executed once per pick-and-place 
cycle on the Raspberry Pi~5 host. The implementation is fully 
deterministic and produces a solution in less than one millisecond, 
which is well below the time required for the physical motors to 
move into position. The remainder of this section describes the 
choice of solution method (Section~8.7.1), the mathematical 
derivation of the solver (Section~8.7.2), the empirical sag 
compensation applied as post-processing (Section~8.7.3), the 
joint limit and reachability constraints (Section~8.7.4), and 
the forward kinematics used for validation and limp-mode 
operation (Section~8.7.5).

\hmsubsubsection{Choice of Solution Method}{UMA}{}

Inverse kinematics for a serial robotic arm can be solved using one 
of three general approaches: a closed-form geometric solution 
specific to the arm's structure, an iterative numerical solution 
based on the manipulator's Jacobian matrix, or a learned solution 
trained from data using a neural network or similar machine 
learning model \cite{craig2005}. The three approaches differ 
significantly in computational cost, predictability, and the 
amount of development effort required to apply them to a given arm.

A closed-form geometric solution was selected for the present 
project. The reasoning behind this choice is summarised in 
Table~6.

\begin{table}[H]
    \centering
    \renewcommand{\arraystretch}{1.3}
    \begin{tabular}{|p{3.0cm}|p{3.5cm}|p{3.5cm}|p{3.5cm}|}
        \hline
        \textbf{Criterion} & \textbf{Geometric (closed-form)} & 
        \textbf{Numerical (Jacobian)} & \textbf{Learned (ML)} \\
        \hline
        Solution type & 
        Exact, single evaluation & 
        Iterative, may not converge & 
        Approximate, depends on training data \\
        \hline
        Compute time & 
        $<$1~ms & 
        1--50~ms per solve & 
        Fast inference, slow training \\
        \hline
        Predictability & 
        Deterministic & 
        Initial-guess sensitive & 
        Non-deterministic at boundaries \\
        \hline
        Implementation & 
        Specific to arm geometry & 
        General, requires tuning & 
        Requires training pipeline \\
        \hline
        Suitability & 
        Ideal for low-DOF arms & 
        Needed for $\geq$6-DOF arms & 
        Useful when geometry is unknown \\
        \hline
    \end{tabular}
    \caption{Comparison of inverse kinematics solution methods. The 
    closed-form geometric approach was selected for this project.}
\end{table}

The selection of the geometric approach is justified by three 
properties of the arm. First, the arm has only four degrees of 
freedom plus the gripper, which is below the threshold at which 
numerical methods become necessary. Standard textbook approaches 
exist for arms of this configuration, requiring only the law of 
cosines and the \texttt{atan2} function from elementary 
trigonometry \cite{craig2005}. Second, the gripper is constrained 
to point directly downwards at all times during pickup operations, 
which removes one degree of freedom from the problem and allows 
the arm's geometry to decompose cleanly into an independent base 
rotation and a planar two-link sub-chain. Third, the geometric 
solution is fully deterministic, producing the same output for 
the same input regardless of execution context, which simplifies 
debugging and testing.

The numerical and learned approaches were considered but rejected. 
The Jacobian-based numerical approach was rejected as unnecessary 
for the four-degree-of-freedom case, since it would introduce 
convergence tuning, singularity handling, and initial-guess 
sensitivity without producing a more accurate solution than the 
closed-form method. The learned approach was rejected because it 
would require collecting thousands of position-angle pairs from 
the physical arm to achieve adequate accuracy, and because its 
behaviour near workspace boundaries would be difficult to predict 
or verify.

\hmsubsubsection{Solution Derivation}{UMA}{}

The arm's kinematic structure is shown in Figure~7. The base joint 
(motor~1) rotates the entire arm in a horizontal plane, while the 
shoulder, elbow, and wrist joints (motors~2, 3, and~4) operate in 
a vertical plane defined by the base angle. The gripper (motor~5) 
opens and closes independently and does not affect the position of 
the gripper tip. The link lengths are 
$L_1 = 25.5$~cm (shoulder to elbow), 
$L_2 = 23.0$~cm (elbow to wrist), and 
$L_3 = 22.0$~cm (wrist to gripper tip), as established in 
Section~7.2.9.

\begin{figure}[H]
    \centering
    \begin{tikzpicture}[
        scale=1.0,
        joint/.style={
            circle, draw, thick, fill=orange!30, 
            minimum size=6mm, inner sep=0pt, font=\small
        },
        link/.style={
            line width=1.2mm, draw=blue!60!black
        },
        target/.style={
            circle, draw, thick, fill=red!50, 
            minimum size=4mm, inner sep=0pt
        },
        dim/.style={
            <->, >=Latex, thin, dashed, gray
        },
        dimlabel/.style={font=\scriptsize\itshape, text=gray},
        anglelabel/.style={font=\scriptsize, text=black!80}
    ]
    
    % Ground line
    \draw[thick, gray!50] (-1, 0) -- (8, 0);
    \draw[pattern=north east lines, pattern color=gray!50, draw=none] 
        (-1, -0.2) rectangle (8, 0);
    
    % Shoulder joint (origin)
    \node[joint] (shoulder) at (0, 2.5) {};
    \node[anglelabel, left=2pt of shoulder] {Shoulder};
    
    % Elbow position (after L1 at angle theta_shoulder)
    \node[joint] (elbow) at (3.5, 4.0) {};
    \node[anglelabel, above right=1pt and 2pt of elbow] {Elbow};
    
    % Wrist position (after L2)
    \node[joint] (wrist) at (5.5, 2.5) {};
    \node[anglelabel, above right=1pt and 2pt of wrist] {Wrist};
    
    % Gripper tip (L3 pointing straight down)
    \node[target] (tip) at (5.5, 0.5) {};
    \node[anglelabel, right=2pt of tip] {Gripper tip};
    
    % Links
    \draw[link] (shoulder) -- (elbow) node[midway, above left, anglelabel] {$L_1$};
    \draw[link] (elbow) -- (wrist) node[midway, above right, anglelabel] {$L_2$};
    \draw[link] (wrist) -- (tip) node[midway, right, anglelabel] {$L_3$};
    
    % Vertical reference at shoulder
    \draw[dashed, gray] (0, 2.5) -- (0, 0);
    \node[dimlabel, left] at (0, 1.25) {$h$};
    
    % Horizontal reference for shoulder angle
    \draw[dashed, gray] (shoulder) -- (1.5, 2.5);
    
    % Shoulder angle arc
    \draw[->, thick, black!70] (0.7, 2.5) arc (0:23:0.7);
    \node[anglelabel] at (1.0, 2.85) {$\theta_2$};
    
    % Elbow angle arc (interior angle)
    \draw[dashed, gray] (elbow) -- ($(elbow)+(0.6,0.25)$);
    \draw[->, thick, black!70] ($(elbow)+(0.5,0.2)$) arc (22:-160:0.5);
    \node[anglelabel] at ($(elbow)+(0.0,0.55)$) {$\theta_3$};
    
    % Target distance d (from shoulder to wrist)
    \draw[dim] (shoulder) -- (wrist);
    \node[dimlabel] at ($(shoulder)!0.5!(wrist) + (0.0,-0.3)$) {$d$};
    
    % Horizontal reach r (from shoulder vertical to wrist)
    \draw[dim] (0, 1.5) -- (5.5, 1.5);
    \node[dimlabel] at (2.75, 1.7) {$r$};
    
    % Z-coordinate of target (relative to shoulder)
    \draw[dim] (6.5, 2.5) -- (6.5, 0.5);
    \node[dimlabel] at (6.8, 1.5) {$z$};
    
    \end{tikzpicture}
    \caption{Geometric model of the arm in the vertical plane after 
    base rotation. The shoulder joint is the origin of the 
    coordinate frame used by the solver. The two-link sub-chain 
    formed by $L_1$ and $L_2$ reaches the wrist position, and the 
    gripper tip is offset by $L_3$ pointing directly downwards. 
    The triangle formed by $L_1$, $L_2$, and the wrist distance 
    $d$ is solved using the law of cosines.}
\end{figure}

The solver computes joint angles in five sequential steps, 
described below.

\textbf{Step 1: Base angle.} The base motor rotates the arm so 
that the vertical plane containing the shoulder, elbow, and 
wrist passes through the target point. Given a target position 
$(x, y, z)$ in the workspace coordinate frame, the base angle 
is computed directly as

\begin{equation}
    \theta_1 = \mathrm{atan2}(y, x)
\end{equation}

This decouples the base rotation from the remainder of the 
problem. The horizontal reach to the target in the rotated 
plane is given by 
$r_{\mathrm{target}} = \sqrt{x^2 + y^2}$.

\textbf{Step 2: End-effector offset.} The gripper points 
directly downwards with the third link $L_3$ aligned along the 
negative $z$-axis. To position the gripper tip at the target 
$(r_{\mathrm{target}}, z)$ in the planar coordinate frame, the 
two-link sub-chain $L_1, L_2$ must place the wrist at

\begin{equation}
    r_{\mathrm{wrist}} = r_{\mathrm{target}}, \qquad 
    z_{\mathrm{wrist}} = z + L_3
\end{equation}

The wrist motor (motor~4) then rotates the gripper to point 
straight down regardless of the angles of motors~2 and~3. The 
$z_{\mathrm{wrist}}$ value is subsequently shifted by the 
shoulder height $h = 35.0$~cm so that the planar inverse 
kinematics is solved with the shoulder joint as the origin.

\textbf{Step 3: Shoulder and elbow angles.} The two-link 
sub-chain forms a triangle with sides $L_1$, $L_2$, and 
$d = \sqrt{r_{\mathrm{wrist}}^2 + z_{\mathrm{wrist}}^2}$, where 
$d$ is the planar distance from the shoulder to the wrist 
position. The interior angle of the triangle at the elbow is 
obtained from the law of cosines:

\begin{equation}
    \cos\theta_e = \frac{L_1^2 + L_2^2 - d^2}{2 L_1 L_2}
\end{equation}

The elbow joint angle is then $\theta_3 = \pi - \theta_e$, 
following the convention that $\theta_3 = 0$ corresponds to a 
fully folded elbow and $\theta_3 = \pi$ corresponds to a fully 
extended arm.

The shoulder angle is obtained by combining two angles. The 
first is the elevation angle $\varphi$ of the line from the 
shoulder to the wrist, measured from the horizontal:

\begin{equation}
    \varphi = \mathrm{atan2}(z_{\mathrm{wrist}}, r_{\mathrm{wrist}})
\end{equation}

The second is the angle $\alpha$ between the shoulder-to-wrist 
line and the upper link $L_1$, also from the law of cosines:

\begin{equation}
    \cos\alpha = \frac{L_1^2 + d^2 - L_2^2}{2 L_1 d}
\end{equation}

The shoulder angle is then $\theta_2 = \varphi + \alpha$.

\textbf{Step 4: Wrist angle.} The wrist must compensate for the 
combined orientation of motors~2 and~3 so that the gripper points 
directly downwards. The required wrist angle is

\begin{equation}
    \theta_4 = -\frac{\pi}{2} - (\theta_2 - \theta_3)
\end{equation}

This expression ensures that the third link $L_3$ remains aligned 
with the negative $z$-axis regardless of the configuration of the 
upper arm.

\textbf{Step 5: Conversion to Dynamixel steps.} Each Dynamixel 
servo accepts a goal position as an integer in the range 
$[0, 4095]$, where the value $2048$ corresponds to the motor's 
neutral position. The conversion from a joint angle in radians 
to a goal position in steps is

\begin{equation}
    \mathrm{steps} = c + \frac{\theta}{\Delta_{\mathrm{rad}}}
\end{equation}

where $c$ is the centre value for the motor in question and 
$\Delta_{\mathrm{rad}} = 2\pi / 4096$ is the angular resolution 
of the encoder. The centre value is $2048$ for motors~1, 2, 4, 
and~5, and $1980$ for motor~3, which has been calibrated to 
account for a small mechanical offset between the encoder zero 
and the physically straight position of the elbow.

The result of the five steps is a dictionary 
$\{m_1, m_2, m_3, m_4, m_5\}$ of integer step values, which is 
serialised by the host and transmitted to the firmware over the 
USB serial connection described in Section~8.5. The gripper 
position $m_5$ is not computed by the solver but is set to 
either the open or closed value depending on the current state 
of the pick-and-place cycle.


\hmsubsubsection{Sag Compensation}{UMA}{}

The geometric solution described in Section~8.7.2 assumes that the 
arm is perfectly rigid, so that the gripper tip reaches exactly 
the position predicted by the joint angles. In practice, the 
physical arm exhibits a measurable amount of gravitational droop, 
or sag, when the arm is extended horizontally under its own 
weight. The sag arises from elastic deformation in the 
3D-printed brackets, small amounts of backlash in the motor 
gearboxes, and mechanical compliance in the daisy-chained motor 
joints. The result is that for a target at horizontal reach $r$, 
the gripper tip arrives at a position that is consistently lower 
than commanded, with the magnitude of the error increasing with 
reach distance.

To compensate for this, an empirical sag model is added as a 
post-processing step. Before the wrist position is passed to the 
two-link inverse kinematics, the commanded $z_{\mathrm{wrist}}$ 
value is increased by a correction that depends on the horizontal 
reach. The correction lifts the target upwards by approximately 
the amount the arm is expected to droop, so that the gripper 
arrives at the originally commanded height after the deformation 
has occurred.

\textbf{Calibration procedure.} The sag model is calibrated by
commanding the arm to a fixed height above the workspace surface
at six different horizontal reach distances between 6 and
36~cm, with sag compensation disabled. The calibration run
that produced the present coefficients used a commanded height
of $z = 2.0$~cm, which caused the gripper to contact the
surface at the three longest reaches. The calibration script's
default height has since been raised to $z = 10.0$~cm to avoid
desk contact during future recalibrations. At each
reach distance, the actual height of the gripper tip is measured 
with a ruler and recorded. The difference between the commanded 
height and the measured height is the empirical droop at that 
reach distance. The full calibration procedure is described in 
Section~8.8.1.

The measurements obtained during calibration of the present arm
are summarised in Table~7. Three reach distances (24, 30, and
36~cm) caused the gripper to contact the surface during
measurement. For the present calibration run, the true droop
values at these distances were extrapolated from the linear
trend established by the three shorter reaches and substituted
into the data set before fitting. The current version of the
calibration script handles desk-hit points differently,
including them directly in the regression fit rather than
replacing them with extrapolated values.

\begin{table}[H]
    \centering
    \renewcommand{\arraystretch}{1.3}
    \begin{tabular}{|c|c|c|c|}
        \hline
        \textbf{Reach (cm)} & 
        \textbf{Commanded $z$ (cm)} & 
        \textbf{Measured $z$ (cm)} & 
        \textbf{Droop (cm)} \\
        \hline
        6.0  & 2.0 & 14.2 & $-12.2$ \\
        \hline
        12.0 & 2.0 & 9.0  & $-7.0$  \\
        \hline
        18.0 & 2.0 & 3.5  & $-1.5$  \\
        \hline
        24.0 & 2.0 & $-2.2$\textsuperscript{*} & $4.2$ \\
        \hline
        30.0 & 2.0 & $-7.5$\textsuperscript{*} & $9.5$ \\
        \hline
        36.0 & 2.0 & $-12.8$\textsuperscript{*} & $14.8$ \\
        \hline
    \end{tabular}
    \caption{Sag calibration measurements. Asterisks denote 
    extrapolated values where the gripper contacted the surface 
    during measurement. The sign convention follows the workspace 
    coordinate frame, where positive $z$ points upwards from the 
    workspace surface.}
\end{table}

\textbf{Model fitting.} Two regression models were fitted to the
measurement data: a linear model
$\Delta z = a r + c$ and a quadratic model
$\Delta z = a r^2 + b r + c$. The calibration script selects
the quadratic model automatically when its root mean square
error is at least 20\% lower than the linear model's error.
For the present calibration data the improvement was marginal
($0.114$~cm versus $0.115$~cm), which does not meet the 20\%
threshold; the quadratic model was nevertheless retained
manually because it captured the slight non-linearity of the
droop at extended reach without increasing the risk of
overfitting. The fitted quadratic model is

\begin{equation}
    \Delta z = -1.49 \times 10^{-4} \, r^2 + 0.912 \, r - 17.77
\end{equation}

where $\Delta z$ is the upward correction in centimetres and 
$r$ is the horizontal reach in centimetres. Figure~8 shows the 
calibration data overlaid with the fitted curve.

\begin{figure}[H]
    \centering
    \begin{tikzpicture}[
        scale=0.95,
        axis/.style={->, >=Latex, thick, black},
        gridline/.style={very thin, gray!30},
        fitcurve/.style={line width=0.8mm, draw=blue!60!black, smooth},
        measured/.style={
            circle, draw, fill=orange!60, thick, 
            minimum size=4mm, inner sep=0pt
        },
        extrapolated/.style={
            circle, draw, fill=white, thick, 
            minimum size=4mm, inner sep=0pt
        },
        axislabel/.style={font=\small},
        ticklabel/.style={font=\scriptsize}
    ]
    
    \def\xs{0.28}
    \def\ys{0.22}
    
    % Grid
    \foreach \x in {0,5,10,15,20,25,30,35,40}
        \draw[gridline] (\x*\xs, -15*\ys) -- (\x*\xs, 17*\ys);
    \foreach \y in {-15,-10,-5,0,5,10,15}
        \draw[gridline] (0, \y*\ys) -- (40*\xs, \y*\ys);
    
    % Zero line
    \draw[thin, gray!60] (0, 0) -- (40*\xs, 0);
    
    % Axes
    \draw[axis] (0, -16*\ys) -- (0, 18*\ys);
    \draw[axis] (-1*\xs, 0) -- (43*\xs, 0);
    
    % X-axis ticks and labels
    \foreach \x in {0,10,20,30,40} {
        \draw[thick] (\x*\xs, -0.1) -- (\x*\xs, 0.1);
        \node[ticklabel, below] at (\x*\xs, -0.18) {\x};
    }
    \node[axislabel, below=22pt] at (20*\xs, -16*\ys) 
        {Horizontal reach $r$ (cm)};
    
    % Y-axis ticks and labels
    \foreach \y in {-15,-10,-5,5,10,15} {
        \draw[thick] (-0.1, \y*\ys) -- (0.1, \y*\ys);
        \node[ticklabel, left] at (-0.18, \y*\ys) {\y};
    }
    \node[axislabel, rotate=90] at (-1.4, 0) 
        {Correction $\Delta z$ (cm)};
    
    % Quadratic fit curve
    \draw[fitcurve] plot[domain=0:40, samples=60] 
        ({\x*\xs}, {(-0.000149*\x*\x + 0.912*\x - 17.77)*\ys});
    
    % Measured points (filled)
    \node[measured] at (6*\xs,  -12.2*\ys) {};
    \node[measured] at (12*\xs,  -7.0*\ys) {};
    \node[measured] at (18*\xs,  -1.5*\ys) {};
    
    % Extrapolated points (hollow)
    \node[extrapolated] at (24*\xs,   4.2*\ys) {};
    \node[extrapolated] at (30*\xs,   9.5*\ys) {};
    \node[extrapolated] at (36*\xs,  14.8*\ys) {};
    
    % Legend (top-left)
    \begin{scope}[shift={(0.5, 3.0)}]
        \draw[fitcurve] (0, 0) -- (0.55, 0);
        \node[ticklabel, right] at (0.6, 0) 
            {Quadratic fit (RMSE $= 0.114$ cm)};
        
        \node[measured] at (0.275, -0.42) {};
        \node[ticklabel, right] at (0.6, -0.42) 
            {Directly measured};
        
        \node[extrapolated] at (0.275, -0.8) {};
        \node[ticklabel, right] at (0.6, -0.8) 
            {Extrapolated (gripper hit surface)};
    \end{scope}
    
    \end{tikzpicture}
    \caption{Empirical sag calibration data and the fitted 
    quadratic correction model. Three points (6, 12, and 18~cm) 
    were measured directly with a ruler. The remaining three 
    points (24, 30, and 36~cm) were extrapolated, since the 
    gripper contacted the workspace surface during measurement at 
    the commanded height of $z = 2.0$~cm and the true droop 
    exceeded the available headroom. The fitted model is 
    $\Delta z = -1.49 \times 10^{-4} r^2 + 0.912 r - 17.77$, with 
    a root mean square error of $0.114$~cm across all six points.}
\end{figure}

\textbf{Runtime application.} The model coefficients are loaded 
automatically at solver initialisation from the 
\texttt{sag\_calibration.json} file produced by the calibration 
procedure, decoupling the model parameters from the source code 
and allowing recalibration without modification of the solver. 
When the solver is invoked without touch calibration data, the 
correction is computed from the target's horizontal reach and 
added to the wrist's $z_{\mathrm{wrist}}$ value before the 
two-link inverse kinematics is solved. In the deployed 
configuration, however, all pick operations and route waypoints 
use per-point surface heights obtained from the touch 
calibration procedure described in Section~8.8.2. Because these 
heights already account for the arm's real-world droop at each 
calibrated distance, the solver's \texttt{skip\_sag} flag is 
set, which bypasses the polynomial sag model entirely and 
avoids double-compensation. The sag model therefore serves as 
a fallback for deployments in which touch calibration data is 
not available.

To prevent the polynomial model from producing dangerously 
large corrections outside its calibrated range, the correction 
is clamped to a safe interval of $[-15, +20]$~cm. A correction 
that would otherwise fall outside this interval is reported as a 
warning and the corresponding pick is treated as a reach 
violation.

\textbf{Limitations.} The sag model is a phenomenological fit 
rather than a physical model of the arm's mechanical behaviour. 
It assumes that the droop depends only on the horizontal reach 
and not on the base angle, the payload, or the orientation of 
the upper arm. It also does not account for hysteresis, where 
the droop at a given reach may differ depending on whether the 
arm approached that reach from a higher or lower position. These 
limitations are acceptable for the present application, in which 
the payload is light and consistent (a 50~mm ball weighing 
approximately 20~g), but they would need to be revisited for an 
arm subjected to heavier or variable loads. In the current 
deployment, the sag model is fully superseded by the per-point 
surface heights recorded during touch calibration 
(Section~8.8.2), which provide a direct empirical correction at 
each calibrated workspace position. The polynomial sag model 
remains in the codebase as a fallback for configurations 
without touch calibration data, but it is not exercised during 
normal automated operation.

\hmsubsubsection{Joint Limits and Reach Clamping}{UMA}{}

Two further constraints are applied to the solver output before 
the motor commands are transmitted to the firmware. The first 
constraint enforces safe operating ranges for each motor, 
preventing the solver from requesting positions that would drive 
the motors into mechanical hard stops or trigger overload errors. 
The second constraint detects targets that lie outside the arm's 
reachable workspace and either rejects them or clamps them to the 
workspace boundary, depending on the operating mode of the 
solver.

\textbf{Joint limits.} Each Dynamixel servo has a configurable 
range of acceptable goal positions. Although the encoder itself 
supports the full $[0, 4095]$ step range, the physical arm 
cannot safely operate across this entire range because the 
3D-printed brackets and the daisy-chained cabling impose 
mechanical limits on rotation. Driving a motor beyond these 
limits would either physically jam the joint or trigger a 
hardware overload error, which latches the motor and requires a 
12~V power cycle to clear. The safe ranges established 
empirically for the present arm are summarised in Table~8.

\begin{table}[H]
    \centering
    \renewcommand{\arraystretch}{1.3}
    \begin{tabular}{|c|l|c|c|c|}
        \hline
        \textbf{Motor} & \textbf{Joint} & 
        \textbf{Centre (steps)} & 
        \textbf{Min (steps)} & \textbf{Max (steps)} \\
        \hline
        m1 & Base    & 2048 & 0    & 4095 \\
        \hline
        m2 & Shoulder & 2048 & 600  & 3500 \\
        \hline
        m3 & Elbow   & 1980 & 463  & 3363 \\
        \hline
        m4 & Wrist   & 2048 & 500  & 3500 \\
        \hline
        m5 & Claw    & 2048 & 0    & 4095 \\
        \hline
    \end{tabular}
    \caption{Per-motor safe operating ranges. The base and gripper 
    motors are unconstrained because their full encoder range is 
    mechanically reachable. The shoulder, elbow, and wrist motors 
    are constrained to prevent the arm from folding into 
    self-collision configurations or overloading under gravity.}
\end{table}

After the solver computes a goal position for each motor, every 
value is checked against the corresponding range. If a value 
falls outside the range, the solver responds in one of two ways 
depending on its operating mode.

\textbf{Reach clamping.} A target position is unreachable if 
the wrist distance $d$ exceeds the maximum reach 
$L_1 + L_2 = 48.5$~cm. In this case, the position is scaled 
inwards along the line from the shoulder to the target until $d$ 
equals 99\% of the maximum reach. This preserves the base angle 
so that the arm still points at the original target direction, 
but limits the planar reach to a value the two-link sub-chain 
can physically achieve. The 1\% margin is retained to avoid 
configurations in which the elbow is fully extended, which can 
trigger mechanical chatter and high motor currents.

A symmetric constraint applies at short range: a target 
position is also unreachable if $d$ is less than the minimum 
reach $|L_1 - L_2| = 2.5$~cm. Targets inside this minimum reach 
are rejected unconditionally rather than clamped, since they 
typically indicate an error elsewhere in the system, such as a 
homography that mapped a ball to coordinates inside the 
shoulder joint itself.

\textbf{Strict and non-strict modes.} The solver provides two 
operating modes that differ in how they handle violations of 
the joint limits and reach constraints. In the default 
non-strict mode, the solver clamps motor values to the safe 
range and emits a warning, allowing operation to continue with a 
slightly modified pose. This is appropriate during interactive 
calibration and manual operation, where a small deviation from 
the requested pose is preferable to interrupting the workflow.

In strict mode, used for production pick-and-place operations 
and for the prevalidation of bin-placement routes, the solver 
raises an exception with a description of the violation rather 
than silently clamping. This fail-closed behaviour ensures that 
the arm never executes a clamped pose during automated operation, 
where the operator may not be present to detect the deviation. 
A clamped pose places the gripper at a different physical 
position than was requested, which can cause the arm to miss the 
ball or collide with bin walls during placement. The strict mode 
makes such failures explicit at the point of computation rather 
than at the point of physical execution.

\hmsubsubsection{Forward Kinematics and Validation}{UMA}{}

The solver also provides a forward kinematics function, which 
computes the gripper tip's Cartesian position from a given set of 
motor positions. Forward kinematics is the inverse operation of 
the inverse kinematics solver: where the inverse kinematics maps 
a Cartesian target to motor positions, the forward kinematics 
maps motor positions back to Cartesian coordinates. The two 
operations together form a closed loop that can be used to 
verify the correctness of the solver and to support operating 
modes in which the motor positions are read from the physical 
arm rather than computed by the solver.

\textbf{Implementation.} The forward kinematics function inverts 
each step of the inverse kinematics derivation in 
Section~8.7.2. Motor step values are converted back to joint 
angles using the inverse of the conversion in Equation~(15), and 
the resulting joint angles are propagated through the kinematic 
chain to compute the position of the gripper tip in the workspace 
coordinate frame. The shoulder height $h$ is added at the end so 
that the result is expressed in the same frame used by the 
inverse kinematics input.

\textbf{Round-trip validation.} The forward kinematics is used 
to verify the solver during automated testing. For a given 
target position $(x, y, z)$ within the reachable workspace, the 
inverse kinematics produces a set of motor positions, which are 
then passed to the forward kinematics. The result should match 
the original target within numerical precision. Any deviation 
indicates a bug in either the inverse or the forward derivation. 
The test suite for the solver runs this round-trip check across 
a swept grid of target positions and verifies that the maximum 
position error is below a small tolerance.

A second test exploits the symmetry of the arm about the $xz$ 
plane. For any target $(x, +y, z)$ in the right half of the 
workspace, there exists a mirror target $(x, -y, z)$ in the left 
half. The two solutions must produce identical shoulder, elbow, 
and wrist angles, and base angles that sum to $4096$ steps when 
expressed in encoder units. The symmetry test verifies this 
property and catches sign errors in the derivation that 
round-trip testing alone might miss.

\textbf{Use in limp-mode calibration.} The forward kinematics is 
also used during a calibration mode in which the operator 
physically guides the arm to a known workspace position by hand. 
In this mode, the torque on motors~1 through~4 is disabled, 
allowing the operator to move the arm freely. Once the gripper 
is at the desired position, the operator triggers a position 
read, in which the firmware reports the current encoder values 
of all motors back to the host. The forward kinematics then 
converts these encoder values to a workspace position, which is 
recorded as a calibration point. This mode is described in detail 
in Section~8.8 and provides a more accurate alternative to manual 
ruler measurement during the touch calibration procedure.